\section{Introdução}
\label{sec:intro}

O desenvolvimento de software sob demanda envolve a contratação de profissionais com competências técnicas específicas, a negociação de escopo, prazo e orçamento, o acompanhamento das entregas e a necessidade de segurança no fluxo de pagamentos. Nesse contexto, mecanismos claros de acompanhamento, registro e proteção financeira são relevantes para reduzir desacordos entre quem contrata o serviço e quem o executa.

O sistema-base deste trabalho é o \textbf{DevMatch}, uma plataforma de \textit{marketplace} que conecta profissionais de TI freelancers a clientes -- empresas ou pessoas físicas -- que necessitam contratar serviços de desenvolvimento de software. A plataforma organiza o processo de contratação desde a publicação de projetos e o envio de propostas até a formalização do contrato, o acompanhamento de marcos de entrega (\textit{milestones}) e a liberação de pagamentos.

Entre as funcionalidades centrais do DevMatch estão a gestão de projetos e propostas, a negociação entre Cliente e Freelancer, a formalização e o acompanhamento de contratos, o pagamento em custódia (\textit{escrow}), a abertura de disputas e a avaliação bidirecional ao final do contrato. O sistema também prevê \textit{matching} inteligente, que recomenda projetos e profissionais conforme critérios de compatibilidade, e verificação de competências por meio de testes técnicos automatizados.

O objetivo deste trabalho é especificar e modelar o DevMatch sob a perspectiva da Engenharia de Software, mantendo consistência entre requisitos e decisões de projeto. Conforme o cronograma da disciplina, esta entrega parcial apresenta a descrição dos atores, o modelo de casos de uso e as histórias de usuário, a arquitetura do sistema, os diagramas de componentes e de implantação e o diagrama de classes. Os diagramas de sequência do sistema, contratos de operação, diagramas comportamentais adicionais e a modelagem de dados serão desenvolvidos nas entregas subsequentes.

\section{Modelagem de Requisitos}
\label{sec:requisitos}

\subsection{Descrição de Atores}
\label{sec:atores}

Os atores foram identificados a partir dos papéis e responsabilidades definidos para o DevMatch. A Tabela~\ref{tab:atores} sintetiza os quatro perfis humanos que interagem diretamente com a plataforma.

\begin{table*}[!htbp]
\centering
\caption{Atores do sistema DevMatch.}
\label{tab:atores}
\small
\begin{tabularx}{\textwidth}{@{}p{3.0cm}p{4.2cm}X@{}}
\toprule
\textbf{Ator} & \textbf{Objetivo no sistema} & \textbf{Principais responsabilidades e interações} \\
\midrule
Freelancer & Encontrar oportunidades e executar serviços de desenvolvimento de software. & Criar e editar perfil profissional; informar \textit{stack}, portfólio, certificações e valor/hora; buscar projetos; enviar e negociar propostas; acompanhar contratos e marcos; realizar testes técnicos; receber pagamentos liberados e avaliar clientes. \\
Cliente & Contratar profissionais para executar projetos de software. & Publicar projetos com escopo, prazo e orçamento; analisar e aceitar propostas; formalizar contratos; acompanhar e aprovar marcos; autorizar liberações de pagamento; abrir disputas e avaliar freelancers. \\
Administrador do Sistema & Garantir governança, segurança, integridade e rastreabilidade da plataforma. & Gerenciar usuários e permissões; apoiar a verificação de identidade; mediar disputas; monitorar logs, auditorias de acesso e transações relevantes. \\
Analista Financeiro/Suporte & Operar processos financeiros e prestar suporte especializado. & Gerenciar custódia e repasses; emitir relatórios financeiros e de comissões; atender chamados relacionados a pagamentos. \\
\bottomrule
\end{tabularx}
\end{table*}

Cliente e Freelancer participam diretamente do ciclo de contratação. Administrador do Sistema e Analista Financeiro/Suporte exercem funções internas de governança, mediação, auditoria e operação financeira. A separação desses papéis orienta as permissões de acesso e reduz sobreposição de responsabilidades.

\subsection{Modelo de Casos de Uso e Histórias de Usuário}
\label{sec:casosuso}

O modelo de casos de uso representa as principais interações dos atores com o DevMatch. Os casos foram derivados das funcionalidades de gestão de projetos, contratação e acompanhamento, pagamento e segurança financeira, avaliação e reputação, \textit{matching} e verificação de competências.

\begin{table*}[!htbp]
\centering
\caption{Casos de uso identificados para o DevMatch.}
\label{tab:casosuso}
\small
\begin{tabularx}{\textwidth}{@{}p{0.9cm}X p{4.4cm}@{}}
\toprule
\textbf{ID} & \textbf{Caso de uso} & \textbf{Ator principal / atores principais} \\
\midrule
UC01 & Autenticar-se na plataforma & Todos os atores humanos \\
UC02 & Criar e editar perfil profissional & Freelancer \\
UC03 & Buscar projetos publicados & Freelancer \\
UC04 & Enviar e negociar proposta & Freelancer \\
UC05 & Realizar teste de competência & Freelancer \\
UC06 & Publicar projeto & Cliente \\
UC07 & Analisar e aceitar propostas & Cliente \\
UC08 & Formalizar contrato e definir marcos & Cliente, Freelancer \\
UC09 & Acompanhar contrato e marcos de entrega & Cliente, Freelancer \\
UC10 & Aprovar marco e autorizar liberação de pagamento & Cliente \\
UC11 & Comunicar-se durante o projeto & Cliente, Freelancer \\
UC12 & Consultar recomendações de \textit{matching} & Cliente, Freelancer \\
UC13 & Abrir disputa & Cliente, Freelancer \\
UC14 & Avaliar a contraparte ao final do contrato & Cliente, Freelancer \\
UC15 & Gerenciar custódia e repasses & Analista Financeiro/Suporte \\
UC16 & Emitir relatórios financeiros e de comissões & Analista Financeiro/Suporte \\
UC17 & Atender chamado relacionado a pagamento & Analista Financeiro/Suporte \\
UC18 & Gerenciar usuários e permissões & Administrador do Sistema \\
UC19 & Mediar disputa & Administrador do Sistema \\
UC20 & Monitorar logs e auditorias & Administrador do Sistema \\
\bottomrule
\end{tabularx}
\end{table*}

Para preservar a legibilidade, o modelo foi dividido em três visões complementares: Freelancer, Cliente e atores internos. Essa divisão não cria funcionalidades duplicadas; casos compartilhados, como acompanhar contrato, abrir disputa e avaliar a contraparte, representam o mesmo comportamento visto por diferentes atores.

\begin{figure*}[!htbp]
\centering
\diagraminclude[0.95\linewidth]{uc_freelancer.pdf}
\caption{Casos de uso associados ao ator Freelancer.}
\label{fig:uc-freelancer}
\end{figure*}

\begin{figure*}[!htbp]
\centering
\diagraminclude[0.95\linewidth]{uc_cliente.pdf}
\caption{Casos de uso associados ao ator Cliente.}
\label{fig:uc-cliente}
\end{figure*}

\begin{figure*}[!htbp]
\centering
\diagraminclude[0.88\linewidth]{uc_internos.pdf}
\caption{Casos de uso associados aos atores internos Administrador do Sistema e Analista Financeiro/Suporte.}
\label{fig:uc-internos}
\end{figure*}

As Figuras~\ref{fig:uc-freelancer}--\ref{fig:uc-internos} mantêm os mesmos identificadores e nomes da Tabela~\ref{tab:casosuso}. A aprovação de um marco é responsabilidade do Cliente; a gestão da custódia e dos repasses é responsabilidade do Analista Financeiro/Suporte; e a mediação de uma disputa é executada pelo Administrador do Sistema, embora a abertura possa ser iniciada por Cliente ou Freelancer.

As histórias de usuário complementam os casos de uso ao explicitar o valor esperado por cada papel.

\begin{table*}[!htbp]
\centering
\caption{Histórias de usuário do DevMatch.}
\label{tab:historias}
\small
\begin{tabularx}{\textwidth}{@{}p{0.9cm}X@{}}
\toprule
\textbf{UC} & \textbf{História de usuário} \\
\midrule
UC02 & Como Freelancer, quero cadastrar e atualizar meu perfil profissional, para que os clientes conheçam minhas competências, portfólio, certificações e valor/hora. \\
UC03 & Como Freelancer, quero buscar projetos publicados, para que eu encontre oportunidades compatíveis com minhas competências. \\
UC04 & Como Freelancer, quero enviar e negociar uma proposta de prazo e valor, para que eu possa concorrer à contratação de um projeto. \\
UC05 & Como Freelancer, quero realizar testes técnicos, para que minhas competências possam ser verificadas e representadas por selos no perfil. \\
UC06 & Como Cliente, quero publicar um projeto com escopo, prazo, orçamento e \textit{stack} desejada, para que freelancers qualificados possam enviar propostas. \\
UC07 & Como Cliente, quero analisar e aceitar propostas recebidas, para que eu escolha o profissional mais adequado ao projeto. \\
UC08 & Como Cliente ou Freelancer, quero formalizar um contrato com marcos de entrega definidos, para que o trabalho possua escopo e critérios de acompanhamento claros. \\
UC09 & Como Cliente ou Freelancer, quero acompanhar os marcos do contrato, para que eu saiba o andamento e as próximas entregas do projeto. \\
UC10 & Como Cliente, quero aprovar um marco concluído, para que o pagamento correspondente possa ser liberado ao Freelancer. \\
UC11 & Como Cliente ou Freelancer, quero manter a comunicação vinculada ao projeto, para que decisões e alinhamentos fiquem registrados. \\
UC12 & Como Cliente, quero receber recomendações de freelancers compatíveis com meu projeto, para reduzir o esforço de busca por profissionais adequados. \\
UC12 & Como Freelancer, quero receber recomendações de projetos compatíveis com meu perfil, para encontrar oportunidades mais relevantes. \\
UC13 & Como Cliente ou Freelancer, quero abrir uma disputa em caso de desacordo, para que o problema possa ser analisado pela plataforma. \\
UC14 & Como Cliente ou Freelancer, quero avaliar a contraparte ao final do contrato, para que o histórico de reputação reflita experiências anteriores. \\
UC15 & Como Analista Financeiro/Suporte, quero gerenciar a custódia e os repasses, para que os valores sejam movimentados conforme a aprovação dos marcos. \\
UC16 & Como Analista Financeiro/Suporte, quero emitir relatórios financeiros e de comissões, para que as operações possam ser acompanhadas e auditadas. \\
UC17 & Como Analista Financeiro/Suporte, quero atender chamados relacionados a pagamentos, para que problemas financeiros sejam tratados pela equipe responsável. \\
UC18 & Como Administrador do Sistema, quero gerenciar usuários e permissões, para que cada perfil tenha acesso somente às funções autorizadas. \\
UC19 & Como Administrador do Sistema, quero mediar disputas, para que conflitos sejam tratados de forma registrada e imparcial. \\
UC20 & Como Administrador do Sistema, quero monitorar logs e auditorias, para que ações relevantes e transações possam ser rastreadas. \\
\bottomrule
\end{tabularx}
\end{table*}

\section{Modelagem de Projeto}
\label{sec:projeto}

\subsection{Arquitetura do Sistema}
\label{sec:arquitetura}

Para o DevMatch, propõe-se uma arquitetura em camadas com módulos de negócio bem delimitados. A organização separa apresentação, controle de acesso, regras de negócio, persistência/auditoria e integrações externas. Essa divisão reduz o acoplamento entre responsabilidades e permite que módulos críticos, como contratos e pagamentos, evoluam de forma independente.

A arquitetura é organizada nas seguintes responsabilidades: (i) \textbf{Apresentação}, composta pelas interfaces Web e móvel; (ii) \textbf{API e Controle de Acesso}, responsável por autenticação, autorização e roteamento; (iii) \textbf{Serviços de Negócio}, com módulos de Usuários, Projetos/Propostas, Contratos/Marcos, Pagamentos/Escrow, Matching, Avaliação/Reputação, Administração e Suporte; (iv) \textbf{Persistência e Auditoria}, responsável pelo armazenamento dos dados e registros rastreáveis; e (v) \textbf{Integrações Externas}, utilizadas para pagamento, notificações e verificação de identidade.

\begin{figure*}[!htbp]
\centering
\diagraminclude[0.93\linewidth]{arquitetura_sistema.pdf}
\caption{Arquitetura proposta para o sistema DevMatch.}
\label{fig:arquitetura}
\end{figure*}

A Figura~\ref{fig:arquitetura} evidencia a separação entre as responsabilidades e o fluxo principal entre as camadas. A modularização favorece escalabilidade e disponibilidade por permitir replicação da aplicação, reforça segurança e privacidade ao centralizar controles de acesso e sustenta a auditabilidade por manter registros das operações críticas.

\subsection{Diagramas de Componentes e de Implantação}
\label{sec:componentes}

\subsubsection{Diagrama de Componentes}

O diagrama de componentes refina a arquitetura ao representar os principais módulos de software e suas dependências. A interface Web/Móvel consome a API da plataforma, que encaminha as requisições aos componentes de domínio. Os componentes de negócio acessam persistência e auditoria e utilizam adaptadores específicos quando precisam se comunicar com provedores externos.

\begin{figure*}[!htbp]
\centering
\diagraminclude[0.94\linewidth]{componentes.pdf}
\caption{Diagrama de componentes do DevMatch.}
\label{fig:componentes}
\end{figure*}

Na Figura~\ref{fig:componentes}, Projetos/Propostas, Contratos/Marcos, Pagamentos/Escrow e Avaliação/Reputação permanecem separados conceitualmente. Essa divisão mantém correspondência com os grupos funcionais dos requisitos e evita concentrar responsabilidades distintas em um único componente.

\subsubsection{Diagrama de Implantação}

O diagrama de implantação descreve a distribuição dos artefatos em nós de execução. A proposta considera clientes Web e móvel acessando a plataforma por HTTPS, um balanceador de carga, uma camada de API, múltiplas instâncias da aplicação, persistência com réplica, \textit{cache}, barramento de eventos e serviços externos fora do ambiente controlado pelo DevMatch.

\begin{figure*}[!htbp]
\centering
\diagraminclude[0.94\linewidth]{implantacao.pdf}
\caption{Diagrama de implantação do DevMatch.}
\label{fig:implantacao}
\end{figure*}

A Figura~\ref{fig:implantacao} explicita os nós de execução e os principais canais de comunicação. As múltiplas instâncias da aplicação e a réplica do banco de dados são coerentes com os requisitos de disponibilidade, escalabilidade e tolerância a falhas previstos para o sistema.

\subsection{Diagrama de Classes}
\label{sec:classes}

O diagrama de classes apresenta uma visão estrutural resumida do domínio do DevMatch. A classe abstrata \textbf{Usuario} concentra dados comuns e é especializada em \textbf{Freelancer}, \textbf{Cliente}, \textbf{Administrador} e \textbf{AnalistaFinanceiro}. O Cliente publica \textbf{Projeto}; cada Projeto recebe \textbf{Proposta} enviadas por Freelancers; uma proposta aceita pode originar um \textbf{Contrato}; e cada Contrato contém um ou mais \textbf{Marco} de entrega.

O fluxo financeiro é representado por \textbf{ContaEscrow} associada ao Contrato e por \textbf{Pagamento} relacionado aos marcos aprovados. O Contrato também pode possuir \textbf{Avaliacao} e \textbf{Disputa}. A Disputa pode ser mediada por um Administrador, enquanto o AnalistaFinanceiro atua na gestão da custódia e dos repasses.

As multiplicidades apresentadas no diagrama indicam que um Cliente pode publicar vários Projetos; um Projeto pode receber várias Propostas; cada Proposta é enviada por um Freelancer; um Contrato relaciona Cliente e Freelancer no contexto de uma proposta aceita; e um Contrato contém um ou mais Marcos. O modelo central foi mantido enxuto para preservar a legibilidade no corpo do artigo.

\begin{figure*}[!htbp]
\centering
\diagraminclude[0.98\linewidth]{classes_resumido.pdf}
\caption{Diagrama de classes resumido do domínio do DevMatch.}
\label{fig:classes}
\end{figure*}
